\section{Instruction Finetuning}

\subsection{基本方法}

收集 (instruction, output) 对，在预训练模型上继续做监督微调。

\begin{importantbox}{Instruction Finetuning 的效果}
FLAN-T5 在 1800+ 任务上进行 instruction finetuning 后，在未见过的任务上表现显著提升。更大的模型从 finetuning 中获益更多。
\end{importantbox}

\subsection{数据策略与对齐}

- 在数据层面，先做 prompt rewriting 和 canonicalization，减少语义重复，再标注最佳回答与备选回答。
- 用多轮对话模拟，增加上下文依赖；对负样本加入多样化错误类型，提升 robustness。
- 采集指令时，同时保留 user intent metadata（任务类型、复杂度、偏好）以便后续 reward model 进行分层建模。

\begin{knowledgebox}{Instruction 数据的可解释化}
把每条 instruction 附上 metadata，使得训练时可以根据意图（如\textit{creative vs. factual}）动态选用 reward signal，减少因单一 loss 导致的泛化误差。
\end{knowledgebox}

\begin{importantbox}{数据质量胜过模型大小}
虽然更大的模型具备更强的底层能力，但  instruction 微调的数据质量决定了最终助手的可靠性。对抗性字段、多样化指令和人工审查是基本保障。
\end{importantbox}

\subsection{Instruction Finetuning 的局限}

\begin{enumerate}
    \item \textbf{开放式生成没有标准答案}：创意写作等任务无法用唯一正确输出来监督。
    \item \textbf{Token 级别的惩罚不区分错误严重程度}：语言建模损失对所有 token 错误一视同仁。
    \item \textbf{人类生成的答案本身不完美}：标注数据有噪声。
\end{enumerate}

\subsection{本章小结}

Instruction finetuning 是从 LM 到助手的第一步，必须配合数据策略与偏好拆分，才能为 RLHF 或 DPO 提供稳定的初始策略。

